\section*{Chapter \ref{ch:rs:portfolio-construction-ii} --- Portfolio Construction II}

\begin{solution}{exo:rs:portfolio-construction-ii:1}
$\sqrt{0.5^2 \times 0.04 + 0.5^2 \times 0.01} = 11.2\%$. With no correlation the equal \omterm{def:rs:portfolio-construction-ii:erc}{risk contributions} are $w_i\sigma_i^2$ equal after normalisation, so
$w_i \propto 1/\sigma_i$: $1/3$ and $2/3$.
\end{solution}

\begin{solution}{exo:rs:portfolio-construction-ii:2}
Without views the posterior mean is $\Pi = \gamma\Sigma w_m$, and the mean--variance optimum on it (with covariance $\Sigma$) is $(\gamma\Sigma)^{-1}\gamma\Sigma w_m = w_m$;
with the posterior covariance $(1+\tau)\Sigma$ it is $w_m/(1+\tau)$, the market scaled.
\end{solution}

\begin{solution}{exo:rs:portfolio-construction-ii:3}
$1/(1 + \phi \times 0.78)$: 0.56 for the reversal, 0.987 for the post-earnings drift, 0.999 for the persistent drift.
\end{solution}

\begin{solution}{exo:rs:portfolio-construction-ii:4}
The average response is set by the inverse covariance and $\gamma$, which the robust term changes little. The largest responses come from names the optimiser
finds nearly interchangeable, where a small alpha change moves weight from one to the other; the penalty on $\sqrt{w^\top\Omega w}$ makes concentrating in either
costly, so the book holds both and the flip is damped (0.33\% to 0.12\% at most).
\end{solution}

\begin{solution}{exo:rs:portfolio-construction-ii:5}
It minimises variance subject to a constraint on the diversity of weights ($\sum\ln w_i$ above a bound): with no constraint the solution is minimum variance,
with an infinitely tight one it is equal weights, and the \omterm{def:rs:portfolio-construction-ii:erc}{equal risk contribution portfolio} is one point on that path (Maillard, Roncalli and Teïletche).
\end{solution}

\begin{solution}{exo:rs:portfolio-construction-ii:6}
Both trade 56\% of the gap each day, but partial adjustment aims at the Markowitz portfolio, dominated by the reversal that decays in a day: it keeps paying to move
towards a target that will have moved. The \omterm{def:rs:portfolio-construction-ii:aim}{aim portfolio} weights the reversal by 0.56 and the slow signals by nearly one, so the target itself moves less.
\end{solution}

\begin{solution}{exo:rs:portfolio-construction-ii:7}
At $\lambda = 1\,000$: daily Markowitz $-22.62$, aim 3.15, partial 1.08, monthly 2.55. At $\lambda = 10\,000$: the aim rule 1.61, the monthly rebalance $-0.68$, the
others below. The aim rule survives because it adjusts both its target (towards slow signals) and its speed (0.23 of the gap at $\lambda = 1\,000$) to the costs.
\end{solution}

\begin{solution}{exo:rs:portfolio-construction-ii:8}
The difference (0.67 against 0.57) is a quarter of the standard error of a Sharpe ratio over 7.9 years (about 0.4). Switch for reasons that are measured
precisely (stability, concentration, turnover), not on this comparison.
\end{solution}

\begin{solution}{pb:rs:portfolio-construction-ii:1}
\begin{enumerate}
\item Mean--variance 0.029\% and 0.33\%; sample covariance 0.028\% and 0.24\%; robust 0.028\% and 0.12\%; Black--Litterman 0.031\% and 0.11\%.
\item 23.9\%, 23.4\%, 22.4\% and 24.4\% of capital, one way.
\item They ignore the alpha; they cannot use it either.
\item The largest jumps between near-interchangeable names; not the average response to the forecast.
\item $\mu = [(\tau\Sigma)^{-1} + P^\top\Omega^{-1}P]^{-1}[(\tau\Sigma)^{-1}\Pi + P^\top\Omega^{-1}q]$; $\tau$ scales the prior's uncertainty relative to the views'.
\item 18.4\% and 17.8\% volatility; HRP holds at most 5.8\% in a name.
\item It orders the assets by clustering and splits capital by inverse variances of sub-blocks, never inverting the matrix.
\item Leverage-averse investors overpay for risky assets, leaving safer ones with higher risk-adjusted returns for investors who can lever (Asness, Frazzini and
Pedersen).
\item 0.57, 0.77, 0.61, 0.72, 0.65, 0.67, 0.66, 0.65; standard error about 0.4.
\item Stability, concentration and costs differ; performance differences are not detectable.
\item The persistent drift ($\phi = 0.00137$ a day), the post-earnings drift (1/60), the one-day reversal (1).
\item It chases the reversal each day, turning over 2.11 times capital a day at a cost of 56.6\% a year.
\item The slow signals at nearly full weight, the reversal at 0.56; it trades 0.56 of the gap a day.
\item \textbf{Named result.} Per basis point: 0.029\% (mean--variance), 0.028\% (sample and robust), 0.031\% (Black--Litterman), 0 for the risk-based rules; at
$\lambda = 100$ the \omterm{def:rs:portfolio-construction-ii:aim}{aim portfolio} nets a Sharpe ratio of 4.06 against $-0.51$ for daily re-optimisation.
\item At $\lambda = 1\,000$ the aim rule nets 3.15, the others 2.55 and below; at $10\,000$ only the aim rule is positive (1.61).
\item It trades rarely, which is cheap when costs are low, but each monthly jump is large, and quadratic costs punish large trades.
\item Noise that is new each day is a fast signal: the aim rule would down-weight it and gain relative to the others.
\item The aim rule, with decay rates estimated per signal and $\lambda$ estimated from the firm's own fills.
\item From the price impact of the firm's trades (chapter~27): regress cost on trade size relative to volume and volatility.
\item The path: with costs, how a book moves decides more of its net performance than which book it moves to.
\end{enumerate}
\end{solution}

\begin{solution}{iq:rs:portfolio-construction-ii:1}
Mean--variance on raw forecasts gives extreme, unstable books. Black--Litterman starts from the returns that make the market optimal and moves away only as far as
views, weighted by their confidence, justify: a Bayesian posterior, whose optimum is the market plus tilts.
\end{solution}

\begin{solution}{iq:rs:portfolio-construction-ii:2}
Allocate so that each asset class contributes equally to risk, ignoring expected returns (assumed roughly proportional to risk). To reach equity-like returns the
low-risk classes, mostly bonds, must be levered; the argument for it is that leverage-averse investors leave safer assets with higher risk-adjusted returns.
\end{solution}

\begin{solution}{iq:rs:portfolio-construction-ii:3}
It subtracts the worst case of the alpha over an \omterm{def:rs:portfolio-construction-ii:robust}{uncertainty set}, which penalises concentration along uncertain directions and damps flips. Size the set from the
forecast's estimated error (a number of standard errors, or a quantile of the chi-squared distribution), and check stability and realised performance.
\end{solution}

\begin{solution}{iq:rs:portfolio-construction-ii:4}
Cluster the assets by correlation distance, order them, and split capital recursively by the inverse variances of the halves. Use it when the covariance matrix is
ill-conditioned or singular, and when stability matters more than using an alpha.
\end{solution}

\begin{solution}{iq:rs:portfolio-construction-ii:5}
Trade partially towards an \omterm{def:rs:portfolio-construction-ii:aim}{aim portfolio} that down-weights the fast signal by $1/(1 + \phi a/\gamma)$ (Gârleanu and Pedersen): aim in front of the target, trade
part of the way.
\end{solution}

\begin{solution}{iq:rs:portfolio-construction-ii:6}
Equal contributions $w_1(\Sigma w)_1 = w_2(\Sigma w)_2$ with $w_1 + w_2 = 1$: $w_1^2\sigma_1^2 + w_1w_2\rho\sigma_1\sigma_2 = w_2^2\sigma_2^2 + w_1w_2\rho\sigma_1\sigma_2$, so
$w_1\sigma_1 = w_2\sigma_2$ for any $\rho$: inverse-volatility weights, $w_1 = \sigma_2/(\sigma_1 + \sigma_2)$.
\end{solution}
